\section{Q\&A 精选：前沿话题讨论}

讲座 Q\&A 环节涉及了多个重要话题，以下是精华摘录：

\subsection{数据枯竭与合成数据}

\textbf{问题}：随着训练数据被不断消耗，可用数据越来越少，未来何去何从？

讲师的回应要点：
\begin{itemize}
    \item 数据许可协议（licensing deals）将成为常态，如 New York Times 与 AI 公司的版权诉讼正在推动这一进程
    \item 企业正在积极寻找``坐拥数据宝藏''的公司进行收购或合作
    \item 合成数据（synthetic data）可以带来一定的训练提升，但过量使用会导致\textbf{模态坍缩}（mode collapse）——模型输出趋于单调
    \item 小型语言模型（Small Language Models）的研究前景广阔：用更少的数据训练针对特定用途的小模型
\end{itemize}

\subsection{RAG：应对幻觉的利器}

\textbf{问题}：是否有集中化的幻觉数据集可以用来防止未来的幻觉？

讲师的回答非常深刻：直接用幻觉数据训练可能适得其反——模型可能学会生成\textbf{更逼真的}幻觉。

\begin{importantbox}{RAG（Retrieval-Augmented Generation）}
讲师认为目前最有效的幻觉缓解方案是 RAG：
\begin{itemize}
    \item 维护一个外部知识库（数据库、向量存储等），存放可信的事实信息
    \item 教模型在生成回答前先从知识库中检索相关上下文
    \item 这种方法巧妙地\textbf{分离了关注点}：
    \begin{itemize}
        \item LLM 负责它擅长的事——生成流畅、连贯的文本
        \item 数据库负责它擅长的事——存储、查找、更新事实信息
    \end{itemize}
    \item 额外好处：当事实发生变化时，只需更新数据库，无需重新训练模型
\end{itemize}
\end{importantbox}

\subsection{LLM 能否推动科学发现？}

\textbf{问题}：随着上下文窗口的指数级增长，LLM 能否通过整合碎片化的研究成果来推动科学突破（如治愈癌症）？

讲师分享了两个视角：
\begin{itemize}
    \item \textbf{直接方式}：让模型直接处理大量论文，寻找跨领域的联系和模式
    \item \textbf{间接方式}：帮助研究者高效消化海量文献（如摘要上千篇论文），加速研究节奏
\end{itemize}

讲师特别提到了一个令人振奋的案例：有研究者将语言模型的框架应用于原子运动模拟（而非自然语言），训练出的``原子语言模型''在从未见过食盐的情况下，给定钠原子和氯原子，正确预测出了\textbf{食盐晶体的结构}。这说明 Transformer 架构的预测能力远不止于语言领域。

\subsection{如何教 LLM 遵守规则？}

\textbf{问题}：如何让 LLM 遵守特定规则（如象棋的合法着法）？

讲师给出了两层方案：

\begin{enumerate}
    \item \textbf{模型层面}：在 fine-tuning 时引入自定义惩罚函数，当模型生成违规输出时施加额外的 loss
    \item \textbf{系统层面}：在模型外部搭建\textbf{预测层 + 策略层}的双层架构。预测层（LLM）负责生成候选输出，策略层（rule engine）负责校验输出是否合规。如果不合规，策略层拒绝并要求 LLM 重新生成
\end{enumerate}

\begin{knowledgebox}{生产级 LLM 系统的标准架构}
讲师指出，几乎所有他见过的生产级 ML 系统——不仅仅是 LLM，甚至在 LLM 之前的传统 ML 系统——都采用了\textbf{预测模型 + 策略层}的双层架构。策略层可以阻止、拒绝或提升（promote）特定的输出，从而将本质上随机的系统（stochastic system）变得可控和可预测。
\end{knowledgebox}

\subsection{本章小结}

Q\&A 环节讨论了当前 LLM 领域的多个前沿话题：数据枯竭推动了合成数据研究和小模型趋势；RAG 通过外部知识库有效缓解幻觉；LLM 的科学应用前景广阔但需要谨慎评估；生产系统中的规则约束需要模型内外的双重保障。

